\documentclass[a5paper, 10pt]{memoir}

\usepackage{color}
\usepackage[english]{babel}
\usepackage[utf8]{inputenc}
\usepackage{amsmath}
\usepackage{graphicx}
\usepackage[lining]{ebgaramond}
\usepackage{makeidx}
\usepackage{fancyvrb}
\usepackage[backend=biber,style=authoryear,citestyle=numeric-comp,sorting=none]{biblatex}\addbibresource{main.bib}

\usepackage[textsize=small, linecolor=magenta, bordercolor=magenta,
            backgroundcolor=magenta, colorinlistoftodos]{todonotes}
\usepackage{letltxmacro}
\LetLtxMacro{\oldtodo}{\todo}
\renewcommand{\todo}[1]{{\color{white}\oldtodo{\textsf{#1}}}}


%\setlength{\parindent}{0cm} % egonw, do we need to have a word about this?
\RecustomVerbatimEnvironment{Verbatim}{Verbatim}{xleftmargin=5mm}


\makeindex

\hyphenation{Java-Script}

\title{A lot of \\ Bioclipse Scripting Language \\ examples \\ (DRAFT)}

\author{E.L. Willighagen, J. Alvarsson}

\begin{document}
\maketitle

\newpage

\mbox{}\vfill
{\noindent \textsf{This work is made available under the Creative Commons Attribution-ShareAlike 4.0 International license.}}
\vskip 5 ex
\begin{center}
	\includegraphics[width=0.3\textwidth]{ccbysa.png}
\end{center}
\vfill
\newpage

\tableofcontents

\chapter{Introduction}
\begin{refsection}

Just to set the record straight, this book is available under a
Creative Commons~4.0 Attribution ShareAlike license. Therefore,
feel free to copy this book, share it with friends, or extend it
with additional code examples.

The next two sections will explain some basics of the Bioclipse
Scripting Language (\textsc{bsl}) idea~\cite{spjuth2009bioclipse}, and the following chapters describe\index{bsl}
functionality available in this language. Each chapter is written
around a particular area.

\section{Scripting Languages}

Bioclipse supports scripting in two programming languages:
JavaScript\index{JavaScript} and Groovy\index{Groovy}.
Support for other programming languages
is possible, but needs some additional funding for development.
Both these languages are enriched with domain extensions.

\section{Extensions}

Extensions of the scripting language are called \emph{managers}.
These managers introduce domain specific functionality, for
example, in the field of cheminformatics.

Besides managers, Bioclipse also extends the scripting language with a few
basic, helpful commands. For example, to get documentation you can
use the \texttt{man} command, for example, for itself:\index{man}
\begin{Verbatim}
man man
\end{Verbatim}
BTW, help is a synonym:\index{help}
\begin{Verbatim}
help man
\end{Verbatim}

Another interesting command is the doi command which opens a web page
with a paper describing the (tool behind the) functionality of a manager.
For example, to get help for the cdk manager, you type:\index{doi}

\begin{Verbatim}
doi cdk
\end{Verbatim}

\section{Domain Objects}

But script extensions are not enough. If they cannot communicate which
each other, then it does not make sense. The primary integration is
dealt with the scripting language.\todo{This is very unclear} But the other thing is that the
managers need a common data model. This is what the domain objects
are about\index{domain object}.

For example, Bioclipse knows the following domain objects:
\vspace{0.25cm}

\begin{tabular}{ l l }
  IMolecule\index{IMolecule}\index{molecule} & a chemical structure \\
  IProtein\index{IProtein}\index{protein} & a protein sequence \\
  IProteinCrystal\index{IProteinCrystal} & a protein crystal structure \\
\end{tabular}



\printbibliography[heading=subbibliography]
\end{refsection}




\chapter{Core functionality}
\begin{refsection}

\section{The bioclipse manager}\index{manager!bioclipse}\index{bioclipse}

One core manager is the \emph{bioclipse} manager. 
For example, it has functionality to convert Eclipse-style paths, based
on projects to full operating system-style paths, which you may
need in the scripting language. For example, given a project in
Bioclipse called ``WikiPathways'' with a folder called ``data'', I
can get the full path with:\index{fullPath}
\begin{Verbatim}
fullPath = bioclipse.fullPath(
  "/WikiPathways/data/" + species + "/"
)
\end{Verbatim}
This full path I can then use with, for example, the Java
\emph{File} class from JavaScript and Groovy. For example,
I am using Groovy in this set up:
\begin{Verbatim}
dataMap = bioclipse.fullPath(
  "/WikiPathways/data/" + species + "/"
);
gpmlFiles = new File(dataMap).listFiles()
\end{Verbatim}
Furthermore, the bioclipse manager has functionality to get some
information about the running Bioclipse. You can get the
version with:\index{version}
\begin{Verbatim}
bioclipse.version()
\end{Verbatim}
And get the location of where the log file is
found:\index{logfileLocation}
\begin{Verbatim}
bioclipse.logfileLocation()
\end{Verbatim}
It can be used to test if there is an Internet
connection:\index{isOnline}
\begin{Verbatim}
bioclipse.isOnline()
\end{Verbatim}
The manager can be used to validate some of these assumptions.
For example, it can be used
to test which Bioclipse version we are running, or make an
assertion on the minimally required version:\index{requireVersion}
\begin{Verbatim}
bioclipse.requireVersion("2.6.1")
\end{Verbatim}
But also fail if Bioclipse does not have active Internet
access:\index{assumeOnline}
\begin{Verbatim}
bioclipse.assumeOnline()
\end{Verbatim}
The bioclipse manager is also the place to look if you would like
to download a web page, such as the front page of the Dutch
newspaper "De Volkskrant":\index{download!webpage}
\begin{Verbatim}
bioclipse.download("http://www.vk.nl/")
\end{Verbatim}


\section{The gist manager}\index{manager!gist}\index{gist}

The gist manager is used for downloading gists. A gist is a
simple way to share small texts on GitHub\index{GitHub},
such as Bioclipse scripts and get verson control on them.
For more information see \url{https://gist.github.com/}.

Originally, a gist was a single file, but that is past
history. The manager followed this change and will now download
all files from a gist:\index{download!gist}

\begin{Verbatim}
gist.download(6282642)
\end{Verbatim}


\section{The ui manager}\index{manager!ui}\index{ui}

The second core manager is the \emph{ui} manager. It has user
interface related functionality. For example, we can use it to
test if a file is available from the Bioclipse workspace:\index{fileExists}

\begin{Verbatim}
knowledgebase = "/WikiPathways/chebi.owl";
if (ui.fileExists(knowledgebase)) {
}
\end{Verbatim}
Importantly, this manager also takes care of interaction with the Bioclipse
user interface itself. For example, we can use it to open editors for
particular blobs of data:\index{open}

\begin{Verbatim}
ui.open(cdk.fromSMILES("c1=cc=cc=c1"))
\end{Verbatim}

\section{The js manager}
\todo{extend}


%\printbibliography[heading=subbibliography]
\end{refsection}




\chapter{Bioinformatics}
\begin{refsection}

\section{The biojava manager}\index{manager!biojava}\index{biojava}

\begin{tabular}{ll}
\textbf{Feature} & Bioinformatics Feature \\
\textbf{Update site} & \url{} \\
\end{tabular}\\
BioJava is a library oriented at \textsc{dna}, \textsc{rna}, and protein
sequences\cite{holland2008biojava,prlic2012biojava}.\index{BioJava}
With this manager we can create data models for sequences, such
as a \textsc{dna} sequence from \textsc{fasta} string:\index{DNAsFromString}

\begin{Verbatim}
dna = biojava.DNAsFromString("> foo\nGCAT")
\end{Verbatim}
Or from just a sequence string, with or without a
name:\index{DNAfromPlainSequence}

\begin{Verbatim}
dna = biojava.DNAfromPlainSequence("GCAT")
dna = biojava.DNAfromPlainSequence(
  "GCAT", "dummy sequence"
)
\end{Verbatim}
With two additional methods, we now have a pipeline to convert
\textsc{dna}\index{DNA} into \textsc{rna}\index{RNA}, and \textsc{rna} into a protein\index{protein} 
sequence:\index{transcriptionOf}\index{translationOf}

\begin{Verbatim}
dna = biojava.DNAfromPlainSequence("GCATATGAA")
rna = biojava.transcriptionOf(dna)
prot = biojava.translationOf(rna)
\end{Verbatim}

\section{The bridgedb manager}\index{manager!bridgedb}\index{bridgedb}

\begin{tabular}{ll}
\textbf{Feature} & Bioclipse BridgeDB \\
\textbf{Update site} & \url{} \\
\end{tabular} \\

BridgeDb is a platform for identifier mapping~\cite{van2010bridgedb}. The
Bioclipse manager makes its functionality available.

At the core, BridgeDb is a framework, but the project
also provides actual identifier mapping databases. And, of course, when you
want to use \textsc{id} mapping functionality, you first need to load such a database.
The plugin is written such that \textsc{id} mapping databases can be downloaded as
Bioclipse plugins, and the extension mechanism allows the manager to list
which mapping databases are available:\index{listIDMapperProviders}
\begin{Verbatim}
dbList = bridgedb.listIDMapperProviders()
\end{Verbatim}

And then the available mapping databases can be loaded, for example,
the first in this example:\index{getIDMapper}

\begin{Verbatim}
mbMapper = bridgedb.getIDMapper(dbList.get(0))
\end{Verbatim}
Mind you, BridgeDb has separate identifier mapping databases for genes and
proteins and for metabolites.

And once we have a mapper then we can start converting identifiers:\index{xref}\index{map}

\begin{Verbatim}
casXref = bridgedb.xref("50-00-0", "Ca")
mappings = bridgedb.map(mbMapper, casXref)
\end{Verbatim}

\printbibliography[heading=subbibliography]
\end{refsection}




\chapter{Cheminformatics}
\begin{refsection}


\section{The cdk manager}\index{manager!cdk}\index{cdk}

\begin{tabular}{ll}
\textbf{Feature} & Bioclipse Chemoinformatics \\
\textbf{Update site} & \url{} \\
\end{tabular}

Basic cheminformatics in Bioclipse is mainly handled by the 
Chemistry Development Kit\index{Chemistry Development Kit}\index{CDK} (\textsc{cdk})~\cite{Steinbeck2003,Steinbeck2006} and for this there 
is the \emph{cdk} manager.

The cdk manager is one with many features. One is to valid\todo{validate?} \textsc{cas}\todo{explain \textsc{cas}} registry number:\index{CAS registry number}

\begin{Verbatim}
cdk.isValidCAS("50-00-0")
\end{Verbatim}
But let's go to the more interesting functionality around chemical graphs.
For example, let's see how we can create molecular structures from a
SMILES string:\index{SMILES}\index{fromSMILES}

\begin{Verbatim}
mol = cdk.fromSMILES("COC")
\end{Verbatim}
Normally, structure diagrams are generated without explicit hydrogens.
But we can easily add them:\index{hydrogens}\index{addExplicitHydrogens}

\begin{Verbatim}
cdk.addExplicitHydrogens(mol)
\end{Verbatim}
We can then calculate a number of properties, including the molecular 
mass\index{molecular mass}, total formal charge\index{charge}, and
molecular formula\index{molecular formula}:\index{calculateMass}\index{totalFormalCharge}\index{molecularFormula}

\begin{Verbatim}
cdk.calculateMass(mol)
cdk.totalFormalCharge(mol)
cdk.molecularFormula(mol)
\end{Verbatim}
Additionally, we can also inspect some of in the information present
in the model:\index{has2d}\index{has3d}\index{isConnected}

\begin{Verbatim}
cdk.has2d(mol)
cdk.has3d(mol)
cdk.isConnected(mol)
\end{Verbatim}
The cdk manager is also central to file support. Before we load it, we
may want to just check the file format\index{file format}:\index{determineFormat}

\begin{Verbatim}
cdk.determineFormat(
  "/ACS Drug Disclosures/AZD5423.cml"
)
\end{Verbatim}
However, this information is not needed when loading files:\index{loadMolecule}

\begin{Verbatim}
mol = cdk.loadMolecule(
  "/ACS Drug Disclosures/AZD5423.cml"
)
\end{Verbatim}
Saving is quite similar, and there are two methods for the two main
formats:

\begin{Verbatim}
cdk.saveCML(mol, "/Test/mol.cml")
cdk.saveMDLMolfile(mol, "/Test/mol.mol")
\end{Verbatim}

\section{The cdx manager}\index{manager!cdx}\index{cdx}


\begin{tabular}{ll}
\textbf{Feature} & Bioclipse CDK Feature \\
\textbf{Update site} & \url{} \\
\end{tabular} \\

The \texttt{cdx} manager is also based on the \textsc{cdk} and exposes functionality
more oriented at \textsc{cdk} developers. For example, we can create a String
representation of the full data model for debugging
purposes:\index{debug}

\begin{Verbatim}
cdx.debug(mol)
\end{Verbatim}
Or we can see the details of the differences between two data
models:\index{diff}

\begin{Verbatim}
cdx.diff(
  cdk.fromSMILES("CC"),
  cdk.fromSMILES("CCC")
)
\end{Verbatim}
And we can list the exact atom types for the atoms in a
molecule:\index{perceiveCDKAtomTypes}\index{atom type}

\begin{Verbatim}
cdx.perceiveCDKAtomTypes(mol)
\end{Verbatim}

\section{The inchi manager}\index{manager!inchi}\index{inchi}

\begin{tabular}{ll}
\textbf{Feature} & Bioclipse Chemoinformatics \\
\textbf{Update site} & \url{} \\
\end{tabular} \\

The \texttt{inchi} manager makes functionality from the InChI
standard\index{InChI}
available~\cite{heller2013inchi,spjuth2013applications}.
The InChI library is not available as a Java library, but is
included as a binary for a selection of platforms and operating
systems. This means that we cannot assume the InChI
functionality is always available in Bioclipse. Furthermore,
we need to load the library:

\begin{Verbatim}
inchi.load()
inchi.isLoaded()
\end{Verbatim}
But when that has succeeded, we can start minting
InChIs:\index{generate!InChI}

\begin{Verbatim}
anInChI = inchi.generate(
  opsin.parseIUPACName("methane")
)
\end{Verbatim}
The returned value is a class called InChI and we can get both
the full InChI as well as the InChIKey from it:\index{InChIKey}

\begin{Verbatim}
fullInChI = anInChI.getValue()
InChIKey = anInChI.getKey()
\end{Verbatim}

\section{The opsin manager}\index{manager!opsin}\index{opsin}

\begin{tabular}{ll}
\textbf{Feature} & Bioclipse Chemoinformatics \\
\textbf{Update site} & \url{} \\
\end{tabular} \\

The opsin manager makes functionality from the \textsc{opsin}\index{OPSIN}
available: convert \textsc{iupac}\index{IUPAC name} names to chemical
structures~\cite{lowe2011chemical}.

\begin{Verbatim}
mol = opsin.parseIUPACName(
  "Ethyl [(1R,3aR,4aR,6R,8aR,9S,9aS)-9-" +
  "{(E)-2-[5-(3-fluorophenyl)-2-pyridinyl]vinyl}-" +
  "1-methyl-3-oxododecahydronaphtho[2,3-c]furan-" +
  "6-yl]carbamate"
)
\end{Verbatim}

\printbibliography[heading=subbibliography]
\end{refsection}






\chapter{Semantic Web}
\begin{refsection}

\section{The rdf manager}\index{manager!rdf}\index{rdf}

\begin{tabular}{ll}
\textbf{Feature} & Resource Description Framework Feature \\
\textbf{Update site} & \url{} \\
\end{tabular} \\

The \texttt{rdf} manager can be used to handle Resource Description
Framework (\textsc{rdf}) data~\cite{willighagen2011linking}.
The basic unit of information in \textsc{rdf} is a
triple and these triples are stored in a \emph{triple store}.
A file based store can be created with this code:\index{createStore}

\begin{Verbatim}
base = rdf.createStore("/tmp/chebiowl")
\end{Verbatim}
Here, the path is a full operating system-style path, and this one
works only on Unix/Linux/BSD systems.

As soon as you have a store, you can start adding triples to it.
For example, when they come from a file, you may want to do
something like this:\index{importFile}

\begin{Verbatim}
knowledgebase = "/WikiPathways/chebi.owl";
kbFormat = "RDF/XML";
base = rdf.createStore("/tmp/chebiowl")
rdf.importFile(base, knowledgebase, kbFormat);
\end{Verbatim}
If the amount of data is limited, you can also create an
in-memory model:\index{createInMemoryStore}

\begin{Verbatim}
knowledgebase = rdf.createInMemoryStore();
\end{Verbatim}
This can be used to create triples from data you are
processing:\index{addObjectProperty}\index{addDataProperty}
\index{object property}\index{data property}

\begin{Verbatim}
rdf.addObjectProperty(knowledgebase,
  "http://linkedchemistry.info/chembl/molecule/m443",
  "http://www.w3.org/2000/01/rdf-schema#subClassOf",
  "http://semanticscience.org/resource/CHEMINF_000000"
)
rdf.addDataProperty(knowledgebase,
  "http://linkedchemistry.info/chembl/molecule/m443",
  "http://www.w3.org/2000/01/rdf-schema#label",
  "CHEMBL268854"
)
\end{Verbatim}

You can get the number of triples held in the store with the
size method:\index{size}

\begin{Verbatim}
rdf.size(knowledgebase);
\end{Verbatim}

And the see the collected triples as Turtle\todo{unclear, explain this a bit more}, you can use:\index{Turtle}\index{asTurtle}

\begin{Verbatim}
rdf.asTurtle(knowledgebase);
\end{Verbatim}


\section{The isbjørn manager}\index{manager!isbjørn}\index{isbjørn}

\begin{tabular}{ll}
\textbf{Feature} & Bioclipse Isbjørn \\
\textbf{Update site} & \url{} \\
\end{tabular} \\

The isbjørn manager uses the Semantic Web to find information about
chemicals, by taking advantage of the Linked Data approaches.

There are two ways to initiate a search: the first is based on a
\textsc{uri} (Uniform Resource Identifier)\index{URI}
for a particular compound, for example, a ChemSpider\index{ChemSpider}
\textsc{uri}:\index{findInfo}

\begin{Verbatim}
knowledgeList = isbjørn.findInfo(
  "http://www.chemspider.com/" +
  "Chemical-Structure.145.rdf#Compound"
)
\end{Verbatim}
But we can also start with a structure:\index{findInfo}

\begin{Verbatim}
knowledgeList = isbjørn.findInfo(
  opsin.parseIUPACName("benzene")
)
\end{Verbatim}
Both methods return a list of \textsc{rdf} stores, one for each data provider.
The results can be saved as a \textsc{html} file with:\index{saveAsHTML}

\begin{Verbatim}
html = isbjørn.saveAsHTML(
  knowledgeList, "/Test/info.html"
)
ui.open(html)
\end{Verbatim}

\section{The owlapi manager}\index{manager!owlapi}\index{owlapi}

\begin{tabular}{ll}
\textbf{Feature} & Resource Description Framework Feature \\
\textbf{Update site} & \url{} \\
\end{tabular} \\

The \texttt{owlapi} manager exposes functionality of the \textsc{owlapi}\index{OWLAPI}
library~\cite{Horridge2011}, to deal with Web Ontology Languages\index{Web Ontology Languages}
(\textsc{owl})\index{OWL} ontologies. The manager is oriented at ontologies
in the Bioclipse workspace. Thus, loading an ontology works
like:\index{load}

\begin{Verbatim}
ontology = owlapi.load(
  "/eNanoMapper/enanomapper.owl", null
);
\end{Verbatim}
The second parameter is a mapper, which can be used to indicate
where imported ontologies can be locally found.
For example:\index{addMapping}

\begin{Verbatim}
mapper = null; // initially no mapper
mapper = owlapi.addMapping(mapper,
  "http://purl.bioontology.org/ontology/npo",
  "/eNanoMapper/npo-asserted.owl"
);
mapper = owlapi.addMapping(mapper,
  "http://www.enanomapper.net/ontologies/" + 
  "external/ontology-metadata-slim.owl",
  "/eNanoMapper/ontology-metadata-slim.owl"
)
ontology = owlapi.load(
  "/eNanoMapper/enanomapper.owl", mapper
);
\end{Verbatim}
You can always see what \textsc{iri}s are mapped with:\index{listMappings}

\begin{Verbatim}
owlapi.listMappings(mapper)
\end{Verbatim}

Once you have loaded an ontology, you can list all
the imported ontologies:\index{getImportedOntologies}

\begin{Verbatim}
imported = owlapi.getImportedOntologies(ontology)
for (var i = 0; i < imported.size(); i++) {
  js.say(
    imported.get(i).getOntologyID().getOntologyIRI()
  )
}
\end{Verbatim}
Similarly, you can list all the classes defined by the
ontology or imported ontologies:\index{showClasses}

\begin{Verbatim}
imported = owlapi.getImportedOntologies(ontology)
for (var i = 0; i < imported.size(); i++) {
  js.say(
    owlapi.showClasses(imported.get(i))
  )
}
\end{Verbatim}
Finally, the \textsc{owlapi} also provide functionality to check for
profile violations:\index{checkVioloations}

\begin{Verbatim}
owlapi.checkVioloations(ontology)
\end{Verbatim}


\printbibliography[heading=subbibliography]
\end{refsection}






\chapter{Online Service Providers}
\begin{refsection}

\section{The \texttt{biows} manager}\index{manager!biows}\index{biows}

\begin{tabular}{ll}
\textbf{Feature} & Bioinformatics Feature \\
\textbf{Update site} & \url{} \\
\end{tabular} \\
The biows manager allows you to retrieve information from biology-related
webservices, including the Uniprot database:\index{UniProt}\index{queryUniProtKB}:

\begin{Verbatim}
biows.queryUniProtKB("P38398")
\end{Verbatim}
The \textsc{embl} nucleotide sequence database at the \textsc{ebi}:\index{queryEMBL}

\begin{Verbatim}
biows.queryEMBL("M10051")
\end{Verbatim}
And also RefSeq:\index{queryRefseq}\index{RefSeq}

\begin{Verbatim}
biows.queryRefseq("NM_000059")
\end{Verbatim}

\section{The chemspider manager}\index{manager!chemspider}\index{chemspider}

\begin{tabular}{ll}
\textbf{Feature} & Bioclipse Chemoinformatics \\
\textbf{Update site} & \url{} \\
\end{tabular} \\

The chemspider manager makes functionality avaiable to interact
with the ChemSpider database\index{ChemSpider}. For example, we
can download a structure with a chemspider identifier
number:\index{download!chemspider}

\begin{Verbatim}
chemspider.download(9606)
\end{Verbatim}
You can also use ChemSpider to resolve InChIKeys to entries
in the database:\index{InChIKey}

\begin{Verbatim}
chemspider.resolve("RCINICONZNJXQF-MZXODVADSA-N")
\end{Verbatim}
Some methods require a security token that can be set via
the preferences (see
\url{http://www.chemspider.com/AboutServices.aspx}).
But when that is done, 
we can search specific compounds:\index{exactSearch}

\begin{Verbatim}
chemspider.exactSearch(
  opsin.parseIUPACName("methane")
)
\end{Verbatim}
Similarly, we can also do a similarity search:

\begin{Verbatim}
chemspider.similaritySearch(
  opsin.parseIUPACName(
    "acetyl salicylic acid"
  ), 0.95
)
\end{Verbatim}
And by substructure:

\begin{Verbatim}
chemspider.substructureSearch(
  opsin.parseIUPACName("benzene")
)
\end{Verbatim}

\section{The opentox manager}\index{manager!opentox}\index{opentox}

\begin{tabular}{ll}
\textbf{Feature} & Bioclipse OpenTox \\
\textbf{Update site} & \url{} \\
\end{tabular} \\

OpenTox is a platform for toxicology\index{OpenTox}, allowing for sharing of data as well
as creating and running computational models\cite{hardy2010collaborative}.
The opentox manager integrates much of its
functionality\cite{willighagen2011computational}.

There are various methods to list what is available. Some information
is available from a central registration server, e.g.
\url{http://apps.ideaconsult.net:8080/ontology/}:\index{listAlgorithms}\index{listDescriptors}\index{listModels}

\begin{Verbatim}
registry = "http://apps.ideaconsult.net:8080/ontology";
algorithms = opentox.listAlgorithms(registry);
descriptors = opentox.listDescriptors(registry);
models = opentox.listModels(registry);
\end{Verbatim}
Similarly, using a particular OpenTox service, e.g. an \textsc{ambit}\index{AMBIT}
instance at
\url{http://apps.ideaconsult.net:8080/ambit2/}\cite{jeliazkova2011ambit},
we can list data sets
and features (properties):\index{listDataSets}\index{listFeatures}

\begin{Verbatim}
service = "http://apps.ideaconsult.net:8080/ambit2/";
dataSets = opentox.listDataSets(service);
features = opentox.listFeatures(service);
\end{Verbatim}
Instead of listing, we can also search for various types, such as
the ToxTree models\index{ToxTree}\cite{patlewicz2008evaluation}:

\begin{Verbatim}
models = opentox.searchModels(registry, "ToxTree");
\end{Verbatim}

We can select one of the available models, and then make a prediction
for a compound:
\begin{Verbatim}
toxTreeModel =
  "http://apps.ideaconsult.net:8080/ambit2/model/3";
opentox.predictWithModel(
  service, toxTreeModel,
  cdk.fromSMILES("CCC")
)
\end{Verbatim}

\section{The pubchem manager}\index{manager!pubchem}\index{pubchem}

\begin{tabular}{ll}
\textbf{Feature} & Bioclipse Chemoinformatics \\
\textbf{Update site} & \url{} \\
\end{tabular} \\

The pubchem manager makes functionality available to interact
with the PubChem database\index{PubChem}. For example, we
can download a structure with a PubChem compound identifier
number:\index{download!pubchem}

\begin{Verbatim}
pubchem.download(2244)
\end{Verbatim}
Or as a 3\textsc{d} structure:\index{download!pubchem 3D}

\begin{Verbatim}
pubchem.download3d(2244)
\end{Verbatim}
And we can search compounds based on a label:\index{search!pubchem}

\begin{Verbatim}
pubchem.search("tamoxifen")
\end{Verbatim}

\section{The openphacts manager}\index{manager!openphacts}\index{openphacts}

Open PHACTS\index{Open PHACTS} is a semantic web-based knowledge platform
under development to
support drug discovery~\cite{Williams2012}. It provides a REST-like interface
of which some of the methods are exposed by the openphacts managers. Because
it is a semantic web platform, compounds, diseases, proteins, and pathways 
are all identified with URIs. Practically, however, we start with a name,
but usign the Open PHACTS ``identity resolution service'' (IRS) we can convert
names into URIs:\index{search!openphacts}

\begin{Verbatim}
cwHits = openphacts.lookUpCompounds("aspirin")
\end{Verbatim}
From the list, we can get a URI with:

\begin{Verbatim}
compoundURI = cwHits.get(0).getURI()
\end{Verbatim}
Another approach to look up a URI for a compound and, as such, find the compound
in the knowledge base, is to start with an IMolecule:

\begin{Verbatim}
compoundURI = openphacts.getURI(
  cdk.fromSMILES("CC(=O)OC1=CC=CC=C1C(=O)O")
)
\end{Verbatim}
If you want to look up information about the search hits, you do:

\begin{Verbatim}
hitsInfo = openphacts.getCompoundsInfo(cwHits)
\end{Verbatim}
Future versions of the manager will likely provide alternative APIs to
get compound information.

Similar compounds can be looked up using the Tanimoto distance measure,
and a similarity cut off of 0.8 by default:

\begin{Verbatim}
similarCompounds = openphacts.findSimilar(
  cdk.fromSMILES("CC(=O)OC1=CC=CC=C1C(=O)O")
)
\end{Verbatim}
If you find that list too long, you can increase the minimal similarity to,
for example, 0.95:

\begin{Verbatim}
similarCompounds = openphacts.findSimilar(
  cdk.fromSMILES("CC(=O)OC1=CC=CC=C1C(=O)O"), 0.95
)
\end{Verbatim}


\printbibliography[heading=subbibliography]
\end{refsection}


\cleardoublepage
\printindex

\end{document}
